%----------------------------------------------------------------------
\section{Discussion}\label{sec:discussion}
%----------------------------------------------------------------------

\subsection{Summary of results}

The following table summarizes the status of ZFC axioms within
the discrimination framework.

\medskip
\begin{center}
\begin{tabular}{l|l|l}
\textbf{ZFC Axiom} & \textbf{Status} & \textbf{Mechanism} \\
\hline
Extensionality & Theorem & Definition of $\kappa$-equivalence \\
Empty set & Theorem & Trivial filtration \\
Pairing & Theorem & Predicate articulation \\
Union & Derived & Requires intersection; $\GF(2)$ gives $\triangle$ \\
Power set & Constructive & Bounded by $\kappa$-state \\
Separation & Unnecessary & Comprehension bounded by algebra \\
Foundation & Partial & Self-membership excluded; cycles $\geq 2$ permitted \\
Infinity & Operationalist axiom (Definition \ref{def:operationalist}) & Indistinguishability (Prop.\ \ref{prop:infinity}) \\
Choice & Constructive (finite); operationalist (infinite) & Boolean dims.\ (Thm.\ \ref{thm:choice}) \\
\end{tabular}
\end{center}

\subsection{Relationship to existing work}

The framework draws on established mathematics:
\begin{itemize}[nosep]
  \item The exterior algebra over $\GF(2)$ and its nilpotency
    properties are standard.
  \item The Yoneda lemma and the view that identity is
    constituted by morphism profiles are classical category theory.
  \item Linear logic's resource sensitivity governs the write
    path; the exponential modality governs the read path.
  \item Homotopy type theory provides the interpretive framework
    for the graded groupoid structure.
\end{itemize}
Additional structural analogies --- to Surreal numbers,
the Grothendieck construction, and CW-complex cell
attachment --- are developed in Appendix \ref{app:analogies}.
These provide interpretive context but are not
load-bearing for any theorem in the paper.

The novelty lies not in any individual component but in the
specific synthesis: discrimination as primitive, valued in
$\GF(2)$ exterior algebra, with Boolean structure as a derived
local property rather than a global assumption, and with an
observer-dependent, monotonically growing regime $\kappa$
that records all discriminations performed.


\subsection{Galois structure of discrimination}
\label{sec:galois}

Every $\kappa$-state splits $V$ into seen and unseen:
\begin{equation}
  0 \to \kappa \to V \to V/\kappa \to 0.
\end{equation}
$\kappa$ is the inner system's regime (what it can
discriminate). $V/\kappa$ is the complement (what it
can't --- the inner-exterior). The exterior algebra
decomposes along this split:
$\bigwedge(V) \cong \bigwedge(\kappa) \otimes
\bigwedge(V/\kappa)$.

The residue filling (Theorem \ref{thm:exhaustive},
Appendix B.18) crosses this tensor boundary:
$R \in \bigwedge(\kappa)$ is a cycle the inner system
can detect but not fill. The filling $R \wedge e_{n+1}$
lives in $\bigwedge(\kappa) \otimes \bigwedge(V/\kappa)$
--- a tensor product straddling the inner/outer boundary.

\emph{The Galois connection.}
Two lattices stand in order-reversing correspondence:
\begin{center}
\begin{tabular}{rcl}
Sub-regimes: & $\kappa_1 \subseteq \kappa_2 \subseteq
  \cdots \subseteq V$ & (more seen) \\
Complements: & $V/\kappa_1 \supseteq V/\kappa_2
  \supseteq \cdots$ & (less unseen)
\end{tabular}
\end{center}
More discriminators $=$ smaller complement $=$ less
freedom in how the unseen could be arranged.
$\kappa$-evolution moves right in the first lattice and
left in the second.

\emph{The Galois group.}
The automorphisms of the complement that preserve the
sub-regime form the Galois group of the extension:
$\mathrm{Gal}(V/\kappa) = GL(V/\kappa,\; \GF(2))$.
This is the symmetry group of what the inner system
\emph{does not know}, holding fixed what it does know.
It acts on the $\gamma = 0$ region (the unseen
discriminators, Section \ref{sec:expanded}).

Each $\kappa$-evolution (adding one discriminator)
\emph{reduces the Galois group}: the stabilizer of the
new discriminator is a proper subgroup. Knowledge
acquisition is Galois-group reduction.

\emph{The Fano automorphism group.}
For the expanded classifier $\GF(2)^3$:
$GL(3, \GF(2)) \cong PSL(2,7)$, the simple group of
order 168, which is the full automorphism group of the
Fano plane $PG(2, \GF(2))$.

The convergence rate of 3 (Theorem \ref{thm:convergence})
is the minimum size of a generating set for $GL(3, \GF(2))$
acting on the Fano plane: three independent interrogations
generate all 168 symmetries of the epistemic structure.
After 3 interrogations, the Galois group is trivial ---
no symmetry of the unknown remains, and the full
structure is determined.

\emph{The Galois correspondence for $\kappa$-states.}
Intermediate sub-regimes $\kappa \subset \kappa' \subset V$
correspond to subgroups of $GL(\dim(V/\kappa), \GF(2))$:
$\kappa = 0$ (nothing seen, full symmetry) $\leftrightarrow$
the full group;
$\kappa = V$ (everything seen, no symmetry)
$\leftrightarrow$ the trivial group. The lattice of
sub-regimes is isomorphic to the lattice of subgroups,
reversed.

This Galois structure was not designed into the
framework. It follows from the short exact sequence of
the $\kappa$-split, which is a property of the exterior
algebra over $\GF(2)$.

\subsection{The bilinear pairing is the one structure}
\label{sec:one-structure}

The meta-analysis of Appendix \ref{app:meta} reveals
that the paper's six closure techniques ---
characteristic 2, Leibniz, linearity, strictness,
the short exact sequence, and the Galois group ---
are six facets of one object: the bilinear evaluation
pairing $\mathrm{ev}(d, a) = \langle d, a \rangle$
over $\GF(2)$. The field ($1{+}1{=}0$), the boundary
compatibility (Leibniz), the second-argument linearity,
the product strictness, the $\kappa$-decomposition,
and the basis-invariance are six views of this pairing.
The paper's 47 audited claims are 47 projections of
one bilinear form onto different parts of the
mathematical landscape.

\subsection{The residue mechanism is uniform}
\label{sec:residue-uniform}

The framework has one mechanism for handling information
that exists at a finer resolution but is invisible at a
coarser one: the \emph{residue}. This mechanism appears
at every level of the paper:

\begin{itemize}[nosep]
  \item \textbf{Set theory} (\S4): gap $(0,0)$ and overlap
    $(1,1)$ in the four-cell profile. A datum whose
    properties are unresolved by the current $\kappa$.
  \item \textbf{Homology} (\S5): $[R] \in H_k$,
    a non-trivial homology class. A $k$-cycle not
    bounding a $(k+1)$-chain.
  \item \textbf{Truncation} (\S5, Cor 5.8): non-unique
    fillings under truncated $\kappa$. Distinct at
    finer resolution, collapsed at coarser.
  \item \textbf{Categories} (\S6): composition data
    lost under 1-categorical projection. The 2-cell
    $d_1 \wedge d_2$ records which discriminations were
    composed; the 1-categorical shadow forgets this.
  \item \textbf{Sheaves} (\S9): presheaf failure of
    the gluing condition. Incompatible local sections
    producing a residue that triggers $\kappa$-evolution.
  \item \textbf{Coverage} (\S9): the $\gamma = 0$ state
    (pre-interrogation). Information that could exist
    but has not been extracted.
  \item \textbf{Meta-level} (\S10): the cofibration is
    the outer tensor product of all domains (sections,
    theorems, appendix verifications), each carrying
    fibers (claims, justifications, dependencies). The
    Cartesian join of these fibers is checked for holes:
    orphaned appendix sections (fibers without claims),
    uncovered spine theorems (claims without fibers).
    This is $\chi = 0$ detection at the meta-level,
    consumed by editorial $\kappa$-evolution.
\end{itemize}

In every case: the residue has $\chi = 0$ (XOR signature
vanishes), it specifies its own filling (the discriminator
or chain that would resolve it), and $\kappa$-evolution
consumes it. The mechanism is the same operation ---
the wedge product detecting non-discrimination ---
applied to different inputs at different grades.

\subsection{Self-audit: the paper as its own implementation}
\label{sec:self-audit}

If the framework models itself, the paper is a
$\kappa$-evolution trace. The spine/appendix structure
instantiates the framework's own write/read distinction:
the spine is the \emph{write path} (each section adds
dimensions, consuming the previous section's residue),
and Appendix \ref{app:verifications} is the \emph{read
path} (each verification section filters the
corresponding spine claim through the algebra, checking
that the profile is well-formed). The cofibration between
spine and appendix is the wedge product
$\text{Spine} \wedge \text{Appendix}$: it identifies
gaps (spine claims without verification), orphans
(verifications without spine claims), and overlaps
(redundant coverage).

\medskip
\noindent\textbf{$\kappa$-evolution trace (spine).}

\begin{center}
\small
\begin{tabular}{p{1.0cm}|p{2.5cm}|p{2.5cm}|p{2.5cm}|p{2.5cm}}
\textbf{Sec.} & \textbf{Triggering residue} &
\textbf{Dimensions added} &
\textbf{(1,0) claims} &
\textbf{Gaps / Overlaps} \\
\hline
\S2 &
Initial: algebra unapplied &
$\GF(2)$, $\bigwedge V$, $v \wedge v = 0$ &
Ground field, nilpotency &
(0,0): no populations \\
\hline
\S3 &
Algebra not applied &
Single-predicate $d$, $\tau/\sigma$ as
datum properties, profile as composed
discriminator, $\kappa$, residue,
concrete model ($V = \GF(2)^3$, $|S| = 4$) &
Four-cell from composition, monoidal read,
higher-level discriminators as dependent types &
(0,0): no standard math \\
\hline
\S4 &
No standard math &
Profile linearity (Definition \ref{def:profile-linear}), operationalist axiom (Definition \ref{def:operationalist}),
extensionality, Russell, foundation, choice &
Russell (via linearity), foundation (partial),
Banach--Tarski &
(0,1): Russell is relocated restriction.
Foundation permits cycles $\geq 2$.
Infinity/choice depend on axiom \\
\hline
\S5 &
Grade structure unused &
$\partial$, $\partial^2 = 0$, $H_k$, exhaustivity &
Chain complex, groupoid &
(1,1): residue $=$ homological hole \\
\hline
\S6 &
Groupoid lacks direction &
Directed morphisms, 2-category, functors,
adjunction, Yoneda, limits &
2-cat axioms, Yoneda &
(0,1): composition is 2-cell, not 1-cell \\
\hline
\S7 &
Direction at grade 1 only &
Triangle identity, toroid, $S_3$, $(n,k)$ &
Free meta-discrimination &
(0,0): dynamics transfer needs proof \\
\hline
\S8 &
Monoidal not connected &
$\otimes$, delooping, CD tower, $[{-},{-}]$ &
Delooping, closure &
(0,0): CD hexagon axioms resolved (recognizer rotation) \\
\hline
\S9 &
Multiple gaps &
$\Omega$, topos, FDE, FOL, self-hosting,
fixed point, sheaves, Fano &
All resolved or pointer to appendix &
(0,0): fixed-point uniqueness resolved ($\dim = 1$) \\
\end{tabular}
\end{center}

\medskip
\noindent\textbf{Cofibration: Spine $\wedge$ Appendix B.}

The cofibration checks that every spine claim (a
$\tau$-assertion) has a matching appendix verification
(a $\sigma$-assertion), and vice versa. The wedge product
$\text{Spine} \wedge \text{Appendix}$ is non-zero iff
there exist claims without verification or verifications
without claims.

\emph{Result}: after iterative $\kappa$-evolution (the
editing process that produced the current draft), the
cofibration is clean:
\begin{itemize}[nosep]
  \item \textbf{0 orphaned appendix sections}. Every one
    of the 23 verification sections in Appendix B is
    referenced by at least one spine section.
  \item \textbf{0 uncovered spine theorems} (in
    \S\S3--6, 8--9). Every theorem has a
    ``(Full verification: \S B.$x$)'' pointer.
  \item \textbf{2 spine theorems with inline-only proofs}
    (\S7: $(n,1)$-structure and free $(n,k)$-structure).
    These are short structural arguments that serve as
    their own verification.
\end{itemize}

The cofibration matrix (spine sections $\times$ appendix
sections) has the following non-zero entries:

\begin{center}
\small
\begin{tabular}{l|l}
\textbf{Spine} & \textbf{Appendix B coverage} \\
\hline
\S3 & Discriminator semantics (B.24),
  concrete model (\S\ref{sec:concrete-model}) \\
\S4 & Russell (B.2), Foundation (B.16), Choice (B.17),
  Banach--Tarski (B.6) \\
\S5 & Exhaustivity (B.18), Groupoid (B.19) \\
\S6 & 2-cat axioms (B.20), Adjunction (B.1),
  Composition (B.3), Yoneda (B.21) \\
\S8 & Delooping/closure (B.22), CD table (B.7) \\
\S9 & Self-enrichment (B.8), Subobj.\ (B.9),
  Power obj.\ (B.4), Proof theory (B.10),
  FOL (B.15), Self-hosting (B.11),
  Fixed point (B.12), Irremovable (B.13),
  Sheaf (B.14), Fano (B.5),
  Boolean/convergence (B.23) \\
\end{tabular}
\end{center}

\medskip
\noindent\textbf{Wedge product of adjacent sections.}

$\S2 \wedge \S3$: clean interface. The algebra
acquires populations. Key semantic decisions made:
discriminators are single predicates, $\tau/\sigma$
are datum properties (not discriminator outputs),
and the membership profile is a composed grade-2
element. These decisions are tracked in
\S\ref{app:discriminator-semantics}.

$\S3 \wedge \S4$: residue consumed (gap/overlap
distinction produces Banach--Tarski).

$\S4 \wedge \S5$: overlap consumed (homology $=$
residues, two concepts identified).

$\S5 \wedge \S6$: residue consumed (groupoid acquires
direction; composition produces 2-cells, acknowledged).

$\S6 \wedge \S7$: clean (1-categories acquire higher
structure via triangle identity).

$\S7 \wedge \S8$: clean (delooping provides interface).

$\S8 \wedge \S9$: residue consumed (monoidal closure
enables self-enrichment; predicate circularity in
fixed-point theorem acknowledged as internal).

\medskip
\noindent\textbf{Identified pendants} (moved to
Appendix \ref{app:analogies}): Surreal framing,
Grothendieck construction, CW-complex interpretation.
Non-load-bearing.

\medskip
\noindent\textbf{Open verifications} (dimensions with
$\gamma = 0$: not yet interrogated):
\begin{itemize}[nosep]
  \item CD hexagon axioms at step 2 (zero divisors may
    kill alternativity early).
  \item CH correspondence with forcing models.
  \item Fixed-point uniqueness (\textbf{resolved}: $\dim = 1$; the $H_2 = 0$
    toroid).
  \item SPPF formalization for implementation.
\end{itemize}

\medskip
\noindent\textbf{Foundational commitments} (the three
axioms, explicitly separated):
\begin{itemize}[nosep]
  \item \textbf{Algebraic}: $\GF(2)$, $\bigwedge V$,
    $x \wedge x = 0$. Mathematical facts.
  \item \textbf{Regulatory}: profile linearity
    (Definition \ref{def:profile-linear}). Structural
    axiom. Excludes Russell, enables bilinear form.
  \item \textbf{Epistemic}: operationalist axiom
    (Definition \ref{def:operationalist}). Foundational
    choice. Governs infinity, infinite choice,
    self-entailment.
\end{itemize}

\medskip
\noindent\textbf{Self-wedge}: $\text{Paper} \wedge
\text{Paper} = 0$. The cofibration analysis found
no contradictions between spine and appendix: every
claim is verified, every verification is claimed.
The paper's $\kappa$-evolution trace, viewed as a
discrimination complex with the spine as $\tau$ and
the appendix as $\sigma$, has trivial self-wedge.

\emph{If a distinction is real, it must be constructible;
if constructible, it must be behaviorally reachable; if
reachable, it must be observable; if observable, it must
be coverable; if not, it is not a valid runtime
distinction.} The paper covers itself.

The self-application operates as three things viewed from
three vertices of the toroid:
\emph{As sheafification}: the paper is a presheaf over
the site of $\kappa$-states (the section-by-section
construction). The cofibration checks the gluing condition.
Previous drafts had structural gaps; the editing process
consumed them. That editing is $\kappa$-evolution.
\emph{As Grothendieck construction}: the SPPF fibers
125 formal objects over the poset of axiom classes.
The axiom-class distribution and dependency chains are
real structural data, not tautologies.
\emph{As fixed-point operator}: self-application yields
the framework. If it yielded something incompatible,
that would be a genuine failure of self-hosting
(Theorem \ref{thm:self-hosting}).
These are one operation from three perspectives.
The self-audit \emph{improved the paper} through
iterative residue consumption. That is the claim worth
making --- not ``the paper passes its own test.''

\begin{remark}[Epistemic limitation of self-audit]
The self-audit detects internal mismatches (orphaned
sections, uncovered theorems, cofibration gaps) but
cannot detect errors of reasoning, equivocations
between levels of abstraction, or insufficient proofs.
A self-consistent framework applied to itself will
always find itself self-consistent. The cofibration
analysis is structural bookkeeping, not mathematical
audit. External review is irreplaceable.
\end{remark}


\subsection{Open question index}
\label{sec:open-questions}

The open questions O1--O6 are stated in the reading guide
(\S\ref{sec:reading-guide}). Here we note their current
status after the full development:

\begin{description}[nosep]
  \item[O1] (\S4, Thm 4.5) \textbf{Resolved.} Scope of evaluation linearity
    exclusion beyond Russell. Partial answer: 8 of 256
    functions admissible for $n = 3$. Full characterization
    of excluded predicate classes: open.
  \item[O2] (\S4, Thm 4.13) \textbf{Resolved.} Useful forms of
    $\in$-induction survive with cycles $\geq 2$.
    Modular $\in$-induction (modulo cycle equivalence):
    open.
  \item[O3] (\S4, Thm 4.31) \textbf{By design.} The operationalist axiom's
    role in infinite choice. The finite case is constructive
    (class \textsf{A}). The infinite case is axiom-dependent
    (class \textsf{AO}). Whether intermediate strengths
    exist: open.
  \item[O4] (\S9, Thm 9.6) \textbf{Resolved.} Join-semilattice property of
    $\mathbf{K}$. Stated as a requirement. Whether
    consistency of profile linearity is preserved under
    arbitrary joins: resolved. $V^*$ is a vector space;
subspace lattice is modular (Appendix \ref{app:risc}).
  \item[O5] (\S9, Thm 9.15) \textbf{Resolved.} Fixed-point
    uniqueness. $\dim \bigwedge^2(\GF(2)^2) = 1$:
    exactly one non-zero grade-2 element exists
    (Appendix \ref{app:fixed-point-stability}).
  \item[O6] (\S8, Thm 8.8) \textbf{Resolved.} Cayley--Dickson hexagon axioms
    at CD step 2. Zero divisors confirmed. Whether
    alternativity survives: open. Whether the Baez--Dolan
    correspondence holds over $\GF(2)^2$: open.
\end{description}

\subsection{Computational considerations}

The exterior algebra $\bigwedge(\GF(2)^n)$ has $2^n$
elements. A naive implementation that materializes the
full algebra would face exponential blowup as $\kappa$
grows. This is a genuine concern, but it rests on a
misunderstanding of what the framework requires at
runtime.

The framework does not require materialization of the
full exterior algebra. It requires:
\begin{enumerate}[label=(\roman*),nosep]
  \item Storage of the discriminators $d_1, \ldots, d_n$
    (the basis of $V$): this is $O(n)$.
  \item Storage of each element's $\tau/\sigma$ profile
    across all dimensions: $O(n)$ per element, $O(n|S|)$
    total.
  \item Computation of wedge products on demand: each
    wedge product of two basis elements is a constant-time
    operation (check linear independence, produce the
    grade-2 element or zero).
\end{enumerate}

No higher-grade element of the exterior algebra needs to
be stored unless it is explicitly queried. The natural
presentation mode is \emph{lazy evaluation}: grade-$k$
elements are computed when needed and discarded when done.
The monoidal read path (Proposition \ref{prop:monoidal})
guarantees that evaluation order does not affect results,
so lazy evaluation is semantically safe.

The total memory required is the raw quantity for storing
the original observations: the population $S$ and the
$\kappa$-regime's discriminators. No additional memory is
required beyond this, because every derived quantity
(residues, homology classes, sheaf sections, Fano
inferences) is computable on demand from the stored
observations.

If memoization or caching of intermediate results is
desired for performance, the natural data structure is
the \emph{shared packed parse forest} (SPPF): a compact
representation of all possible parse trees that shares
common sub-structures. The full formalization of an SPPF
under the discrimination regime --- showing that the SPPF
structure is itself a $\kappa$-equivalence class of parse
trees, and that its sharing structure corresponds to
$\kappa$-deformations that identify isomorphic sub-trees
--- is an open question for implementation work.

%----------------------------------------------------------------------
\section*{Acknowledgments}
%----------------------------------------------------------------------

This paper formalizes ideas developed through extended collaborative
dialogue. The discrimination framework, the identification of
$\GF(2)$ as the ground field, the observer-dependent $\kappa$-evolution
principle, and the philosophical position that Boolean structure
is derived rather than given are due to the originator of the
framework.

